% Original analytic teaching diagram, 2026-11-12.
% One Monte Carlo computation as a pipeline: base uniforms -> generated sample -> integrand value -> average,
% and the three places where variance control intervenes (distribution, formula, sampling mechanism).
\definecolor{VB}{HTML}{7188A3}\definecolor{VBF}{HTML}{E5EAF0}
\definecolor{VG}{HTML}{829C8B}\definecolor{VGF}{HTML}{E8EEE9}
\definecolor{VY}{HTML}{B4A16F}\definecolor{VYF}{HTML}{F1ECDD}
\definecolor{VR}{HTML}{A7787D}\definecolor{VRF}{HTML}{F0E5E6}
\definecolor{VP}{HTML}{9588A5}\definecolor{VPF}{HTML}{ECE8F0}
\definecolor{VI}{HTML}{3B4252}\definecolor{VM}{HTML}{6B7280}\definecolor{VL}{HTML}{4C4D4F}
\begin{tikzpicture}[x=1mm,y=1mm,>=Stealth,font=\figurefont\fontsize{8}{10}\selectfont,
lab/.style={text=VM,font=\figurefont\fontsize{7.5}{9}\selectfont},
box/.style={draw=VB,fill=VBF,line width=.6pt,rounded corners=1.2mm,minimum height=9mm,minimum width=17mm,align=center,text=VI,inner sep=1.2mm},
ctl/.style={draw=VY,fill=VYF,line width=.6pt,rounded corners=1.2mm,minimum height=10mm,text width=24mm,align=center,text=VI,inner sep=1.2mm},
arr/.style={-{Stealth},draw=VL,line width=.7pt},
hit/.style={-{Stealth},draw=VY,line width=.8pt}]
\node[box] (u) at (0,0) {基础随机数\\$\bm{U}$};
\node[box] (x) at (27,0) {样本\\$X=T(\bm{U})$};
\node[box] (f) at (54,0) {函数值\\$f(X)$};
\node[box] (i) at (81,0) {平均\\$\hat{I}_n$};
\draw[arr] (u) -- (x); \draw[arr] (x) -- (f); \draw[arr] (f) -- (i);
\node[ctl] (c1) at (27,-21) {改变分布\\$X\sim q$，乘权重$w$};
\node[ctl] (c2) at (54,-21) {改变计算公式\\控制变量、条件期望};
\node[ctl] (c3) at (0,-21) {改变抽样机制\\分层、点集、成对};
\draw[hit] (c1.north) -- (x.south);
\draw[hit] (c2.north) -- (f.south);
\draw[hit] (c3.north) -- (u.south);
\node[lab,anchor=south] at (40.5,6.5) {固定$I=\mathbb{E}_P[f(X)]$与总预算};
\end{tikzpicture}
